\section*{Capítulo \ref{ch:b2:michelson} --- El interferómetro de Michelson}

\begin{solution}{exo:b2:michelson:1}
$N = 2\Delta e/\lambda = 340$; $250$ franjas: $\Delta e = 125\lambda = \qty{74}{\micro m}$; una décima de
franja son $\lambda/20 = \qty{30}{nm}$.
\end{solution}

\begin{solution}{exo:b2:michelson:2}
$p_0 = 2e/\lambda = 1697.8$: no es entero, el centro no es brillante ($\varepsilon =
0.8$). $r_k = f\sqrt{(k - 1 + \varepsilon)\lambda/e}$: $3.1$, $4.6$, \qty{5.7}{cm}. Con \qty{0.05}{mm} los
radios crecen en $\sqrt{10}$: pocos anillos y anchos.
\end{solution}

\begin{solution}{exo:b2:michelson:3}
$\lambda/2\alpha = \qty{3.2}{mm}$; seis franjas en \qty{2}{cm}; doblar $\alpha$ reduce a la mitad la
separación. Contacto: trasládese el espejo hasta que aparezcan las franjas de luz
blanca (una blanca flanqueada por coloreadas) --- solo existen dentro de
un micrómetro de $\delta = 0$.
\end{solution}

\begin{solution}{exo:b2:michelson:4}
$2(n - 1)e/\lambda = 177$ franjas; $0.2$ franjas dan $\Delta n = 0.2\lambda/2e = \qty{6e-4}{}$.
Un centímetro de vidrio añade milímetros de camino dispersivo: los colores
ya no concuerdan en ninguna parte, no hay franja blanca.
\end{solution}

\begin{solution}{exo:b2:michelson:5}
(a) $\lambda^2/2\Delta\lambda = \qty{0.291}{mm}$. (b) $5.00/17 = \qty{0.294}{mm}$: $\Delta\lambda = \lambda^2/2\Delta e =
\qty{0.590}{nm}$; $0.1\%$: $\pm\qty{0.0006}{nm}$. (c) En $e = \qty{0.145}{mm}$ (la primera
desaparición): $p_1 = 492.4$, $p_2 = 491.9$ --- medio orden de diferencia; en
\qty{0.29}{mm} difieren en uno. (d) Sobreviven las franjas más débiles: $V_{\min} =
(I_1 - I_2)/(I_1 + I_2) = 1/3$.
\end{solution}

\begin{solution}{exo:b2:michelson:6}
(a) $2(n - 1)\ell = 94\lambda$: $n - 1 = \qty{2.77e-4}{}$. (b) $0.94$ franjas por cada uno por ciento;
$\dd(n - 1)/\dd T = -(n - 1)/T = \qty{-9.5e-7}{K^{-1}}$. (c) Unos pocos milibares o un
grado en un brazo mueven las franjas: de ahí el vacío o los brazos iguales. (d)
$0.1/94$: $\pm\qty{3e-7}{}$.
\end{solution}

\begin{solution}{exo:b2:michelson:7}
(a) $\ell_c = \lambda^2/\Delta\lambda = \qty{3}{cm}$: los anillos se apagan más allá de $e \approx \qty{1.5}{cm}$. (b)
\qty{0.3}{mm}, $e \approx \qty{0.15}{mm}$. (c) $\ell_c = \qty{300}{m}$: sí. (d) $\ell_c \approx \qty{1}{\micro m}$:
$2e < \qty{1}{\micro m}$, un par de órdenes.
\end{solution}

\begin{solution}{exo:b2:michelson:8}
(a) $2(n - 1)e_g = \qty{5.2}{mm}$; sí, con \qty{2.6}{mm} de recorrido del espejo. (b) El
camino en el vidrio varía con $\lambda$ en $2e_g\Delta n = \qty{100}{\micro m} \gg \ell_c$: no hay
cero independiente de la longitud de onda. (c) $2(n - 1)\Delta e_g < \qty{0.3}{\micro m}$: $\Delta e_g <
\qty{0.3}{\micro m}$ --- una inclinación relativa por debajo de $\qty{1.5e-5}{rad}$ en \qty{2}{cm}. (d) Una
película tiene micrómetros de espesor; un cubo es simétrico: ambos haces ven
el mismo vidrio.
\end{solution}

\begin{solution}{exo:b2:michelson:9}
(a) $\Delta t \approx Lv^2/c^3 = \qty{4.1e-15}{s}$, un camino de \qty{1.2}{\micro m}. (b) Intercambiar los
brazos lo duplica: $2Lv^2/c^2\lambda = 0.4$ franjas. (c) $v < 30\sqrt{0.01/0.4} =
\qty{4.7}{km/s}$. (d) No se ve ningún movimiento por un éter; la celeridad de la luz
es la misma por ambos brazos --- el postulado de la relatividad.
\end{solution}

\begin{solution}{exo:b2:michelson:10}
(a) $I = S_0[1 + \cos(4\pi\nu_0e/c)]$, de periodo $\lambda/2$ en $e$. (b) Cada frecuencia
da una franja de periodo algo distinto; una banda gaussiana de anchura
$\Delta\nu$ suma franjas con una envolvente gaussiana de anchura $\sim c/2\Delta\nu$.
(c) $R = \nu/\Delta\nu = 2e_{\max}/\lambda = 2 \times 10^5$. (d) Todas las longitudes de onda caen sobre
un solo detector a la vez (Fellgett) por una abertura ancha (Jacquinot):
mucha más señal que un espectrómetro de rendija en el infrarrojo,
pobre en fotones.
\end{solution}

\begin{solution}{exo:b2:michelson:11}
(a) El camino de cada brazo cambia en $2N\Delta L$, en sentidos opuestos: $\Delta\delta = 4N\Delta L
= \qty{4.8e-15}{m}$, $\Delta\varphi = 2\pi\Delta\delta/\lambda = \qty{2.8e-8}{rad}$. (b) Cerca de una franja
oscura el detector ve una intensidad proporcional a $\Delta\varphi^2$ --- o,
con un pequeño desplazamiento, a $\Delta\varphi$: un cambio relativo del orden de $10^{-8}$. (c)
$5 \times 10^{20}$ fotones por segundo, $5 \times 10^{17}$ por milisegundo, con ruido
relativo $1.4 \times 10^{-9}$: la señal es veinte veces el ruido; más potencia
baja el ruido como $1/\sqrt P$. (d) $\Delta L \propto L$; las suspensiones pendulares
filtran el movimiento del suelo por encima de su resonancia.
\end{solution}

\begin{solution}{exo:b2:michelson:12}
(a) Para un punto $P$ a la profundidad $e(P)$ y con incidencia $i$, $\delta = 2e(P)\cos i$
más un término $\propto\alpha i$ debido a la segunda reflexión inclinada, que
se anula con incidencia normal: solo allí llegan al ojo los dos rayos que salen de $P$
--- localización en los espejos. (b) $ei^2 < \lambda/4$: $i_{\max} <
\sqrt{\lambda/4e} = \qty{0.09}{rad}$. (c) Los anillos dependen solo de $i$: cualquier posición
del ojo ve el mismo patrón angular. (d) La cuña vista fuera de los
espejos; enfóquese sobre $M_2$ o redúzcase la inclinación.
\end{solution}

\begin{solution}{pb:b2:michelson:1}
\textbf{1.} Franjas de cuña; $\lambda/2\alpha = \qty{2}{mm}$: $\alpha = \qty{1.6e-4}{rad}$.

\textbf{2.} $M_1'$ se ha vuelto paralelo a $M_2$: la lámina de aire, cuyas
\omterm{prop:b2:two-wave-interference:film}{franjas de igual inclinación} son anillos.

\textbf{3.} $158\lambda/2 = \qty{50}{\micro m}$; anillos tragados por el centro: $e$
disminuye.

\textbf{4.} El contacto, $e = 0$. La luz de sodio no muestra nada especial allí (una
sola longitud de onda, campo uniforme en cualquier caso); la luz blanca señala la única
posición en la que todos los colores tienen $\delta = 0$, flanqueada por colores a medida que los
órdenes se separan.

\textbf{5.} $\delta = 2\alpha x$: una franja blanca sobre la línea de contacto y una o
dos coloreadas a cada lado ($2\alpha x < \ell_c \approx \qty{1}{\micro m}$).

\textbf{6.} $r_{10} = f\sqrt{10\lambda/e} = \qty{5}{mm}$: $e = 10\lambda f^2/r^2 = \qty{1.0}{cm}$.

\textbf{7.} Paralela, los dos haces cruzan el mismo vidrio para toda
longitud de onda; sin $C$, la dispersión de \qty{5}{mm} de vidrio no deja ningún
cero común: no hay franja blanca.

\textbf{8.} $\lambda^2/4\Delta\lambda = \qty{0.145}{mm}$; luego, cada $\lambda^2/2\Delta\lambda = \qty{0.291}{mm}$.

\textbf{9.} $11$ intervalos en \qty{3.20}{mm}: \qty{0.291}{mm}, $\Delta\lambda = \qty{0.597}{nm}$;
$\pm\qty{10}{\micro m}$ sobre \qty{3.2}{mm}: $\pm\qty{0.002}{nm}$.

\textbf{10.} $492.4$ y $491.9$ (medio orden de diferencia); en el retorno,
$984.7$ y $983.7$.

\textbf{11.} $1/3$.

\textbf{12.} $\ell_c = \lambda^2/\Delta\lambda = \qty{1.7}{cm}$: $2e < \qty{1.7}{cm}$, unas $29$.

\textbf{13.} Tres sistemas de franjas: contraste pleno cada $\lambda^2/2\Delta\lambda$, con
dos caídas parciales entre medias --- el patrón de tres rendijas.

\textbf{14.} $2(n - 1)\ell = \qty{29}{\micro m}$: $46$ franjas.

\textbf{15.} $n - 1 = 46\lambda/2\ell = \qty{2.91e-4}{}$.

\textbf{16.} $9.2$ franjas.

\textbf{17.} Con brazos iguales, la deriva es común a ambos y se cancela; una
desigualdad de \qty{1}{cm} da $2 \times 0.01 \times 3 \times 10^{-6} = 0.1$ franjas ---
inocua para ambas partes.

\textbf{18.} Están presentes, sin cambios, durante toda la cuenta: un
desplazamiento constante.

\textbf{19.} $71$ franjas; la cuenta da $n - 1$, una firma del gas.

\textbf{20.} $31\,600$ franjas en un centímetro, leídas hasta una décima:
\qty{30}{nm}, $3 \times 10^{-6}$; el tornillo da un micrómetro.

\textbf{21.} $\delta$ oscila $\pm\qty{100}{nm}$, $\pm0.16$ franjas a \qty{100}{Hz}: el
ojo ve el contraste rebajado; a $25$ imágenes por segundo la cámara
muestrea cuatro periodos por imagen y ve un patrón congelado o que late
despacio.

\textbf{22.} $10^{-8} \times 31\,600 = 3 \times 10^{-4}$ franjas: nada.

\textbf{23.} $\lambda/20 = \qty{32}{nm}$ en \qty{10}{mm}: $3 \times 10^{-6}$.

\textbf{24.} Compara una longitud directamente con una longitud de onda, franja a
franja; la red despliega todas las longitudes de onda a la vez en un intervalo amplio.

\textbf{25.} Longitud: anillos contados en el centro. Doblete: la
pérdida periódica de contraste. Índice: franjas que pasan mientras se llena
la cubeta. \omterm{prop:b2:scalar-light-model:coherence}{Longitud de coherencia}: el espesor para el que se apagan los anillos.
\end{solution}
